% GENERATED FILE. Edit canonical lessons, then run npm run generate:books.
% canonical-source-hash: fnv1a64:37bcb583106f07de
% canonical-lessons: PA-C182-somvar, PA-C182-mangalvar, PA-C182-budhvar, PA-C182-veervar, PA-C182-shukkarvar, PA-C182-shanivar, PA-C182-aitvar, PA-C182-airhaa-ai, PA-C182-weekday-review, PA-C182-payoff

\chapter{A Week on the Noticeboard}
\label{ch:pa-a2-weekday-schedule}
\hlchaptermodality{182}

\begin{chapteropening}
\textbf{By the end of this chapter:} \emph{I can recognise the seven weekdays and understand one simple opening-days notice.}

Retrieve the full week and a two-line opening-days notice through listening, speaking, reading, and bounded Gurmukhi writing after an atom-free review.
\end{chapteropening}

\section[\texorpdfstring{\pa{ਸੋਮਵਾਰ} (somvaar)}{ਸੋਮਵਾਰ (somvaar)}]{\pa{ਸੋਮਵਾਰ} — Monday}

\label{lesson:PA-C182-somvar}



\begin{quote}
\emph{Hear:} \emph{somvaar}; YOU SAY: \textbf{\pa{ਸੋਮਵਾਰ}} Start with the sound. It means \emph{Monday}.
\end{quote}

\begin{grammarlens}[title={Grammar Lens: one weekday}]
\begin{quote}
\textbf{\pa{ਸੋਮਵਾਰ}} — \emph{somvār} — Monday
\end{quote}

Keep it as one calendar label. A full week comes slowly, one day at a time.
\end{grammarlens}

\begin{cousinweb}
\emph{Somvār} joins Sanskrit \emph{soma}, \textquotedblleft{}moon\textquotedblright{}, with \emph{vāra}, \textquotedblleft{}day of the week\textquotedblright{} (CDIAL 11547): the moon day.
\end{cousinweb}

\subsection*{Your turn}
\begin{itemize}
\raggedright
  \item \emph{Hear:} \emph{somvaar}; YOU SAY: \textbf{\pa{ਸੋਮਵਾਰ}}
  \item \emph{Read it:} \textbf{\pa{ਸੋਮਵਾਰ}}
\end{itemize}

\begin{tcolorbox}[breakable,colback=teal!4,colframe=teal!35!black,title={Before you move on}]
Say \emph{Monday}. (\textbf{\pa{ਸੋਮਵਾਰ}}.)

Sources: Punjabi University, \href{https://dic.learnpunjabi.org/default.aspx?look=\%E0\%A8\%B8\%E0\%A9\%8B\%E0\%A8\%AE\%E0\%A8\%B5\%E0\%A8\%BE\%E0\%A8\%B0}{Punjabi-English Dictionary}, for the weekday; Turner, \href{https://dsal.uchicago.edu/dictionaries/soas/}{CDIAL 11547}, for \emph{vāra}.
\end{tcolorbox}
\section[\texorpdfstring{\pa{ਮੰਗਲਵਾਰ} (mangalvaar)}{ਮੰਗਲਵਾਰ (mangalvaar)}]{\pa{ਮੰਗਲਵਾਰ} — Tuesday}

\label{lesson:PA-C182-mangalvar}



\begin{quote}
Retrieve Monday: \textbf{\pa{ਸੋਮਵਾਰ}}.
\end{quote}

\begin{grammarlens}[title={Grammar Lens: one weekday by ear}]
\emph{Hear:} \emph{mangalvaar}; YOU SAY: \textbf{\pa{ਮੰਗਲਵਾਰ}}

\begin{quote}
\textbf{\pa{ਮੰਗਲਵਾਰ}} — \emph{maṅgalvār} — Tuesday
\end{quote}
\end{grammarlens}

\begin{cousinweb}
\emph{Maṅgalvār} joins Sanskrit \emph{maṅgala}, the name of Mars, with \emph{vāra}, \textquotedblleft{}day of the week\textquotedblright{} (CDIAL 11547).
\end{cousinweb}

\subsection*{Your turn}
\begin{itemize}
\raggedright
  \item \emph{Hear:} \emph{mangalvaar}; YOU SAY: \textbf{\pa{ਮੰਗਲਵਾਰ}}
  \item \emph{Read it:} \textbf{\pa{ਮੰਗਲਵਾਰ}}
  \item \emph{Say it:} \textbf{\pa{ਸੋਮਵਾਰ}}, then \textbf{\pa{ਮੰਗਲਵਾਰ}}
\end{itemize}

\begin{tcolorbox}[breakable,colback=teal!4,colframe=teal!35!black,title={Before you move on}]
Say Monday, then Tuesday.

Sources: Punjabi University, \href{https://dic.learnpunjabi.org/default.aspx?look=\%E0\%A8\%AE\%E0\%A9\%B0\%E0\%A8\%97\%E0\%A8\%B2\%E0\%A8\%B5\%E0\%A8\%BE\%E0\%A8\%B0}{Punjabi-English Dictionary}, for the weekday; Turner, \href{https://dsal.uchicago.edu/dictionaries/soas/}{CDIAL 11547}, for \emph{vāra}.
\end{tcolorbox}
\section[\texorpdfstring{\pa{ਬੁੱਧਵਾਰ} (budhvaar)}{ਬੁੱਧਵਾਰ (budhvaar)}]{\pa{ਬੁੱਧਵਾਰ} — Wednesday}

\label{lesson:PA-C182-budhvar}



\begin{quote}
Say \textbf{\pa{ਸੋਮਵਾਰ}}, then \textbf{\pa{ਮੰਗਲਵਾਰ}}.
\end{quote}

\begin{grammarlens}[title={Grammar Lens: one weekday by ear}]
\emph{Hear:} \emph{budhvaar}; YOU SAY: \textbf{\pa{ਬੁੱਧਵਾਰ}}

\begin{quote}
\textbf{\pa{ਬੁੱਧਵਾਰ}} — \emph{budhvār} — Wednesday
\end{quote}
\end{grammarlens}

\begin{cousinweb}
\emph{Budhvār} joins Sanskrit \emph{budha}, the name of Mercury, with \emph{vāra}, \textquotedblleft{}day of the week\textquotedblright{} (CDIAL 11547).
\end{cousinweb}

\subsection*{Your turn}
\begin{itemize}
\raggedright
  \item \emph{Hear:} \emph{budhvaar}; YOU SAY: \textbf{\pa{ਬੁੱਧਵਾਰ}}
  \item \emph{Read it:} \textbf{\pa{ਬੁੱਧਵਾਰ}}
  \item \emph{Say it:} \textbf{\pa{ਸੋਮਵਾਰ}}, \textbf{\pa{ਮੰਗਲਵਾਰ}}, \textbf{\pa{ਬੁੱਧਵਾਰ}}
\end{itemize}

\begin{tcolorbox}[breakable,colback=teal!4,colframe=teal!35!black,title={Before you move on}]
Say \emph{Wednesday}. (\textbf{\pa{ਬੁੱਧਵਾਰ}}.)

Sources: Punjabi University, \href{https://dic.learnpunjabi.org/default.aspx?look=\%E0\%A8\%AC\%E0\%A9\%81\%E0\%A9\%B1\%E0\%A8\%A7\%E0\%A8\%B5\%E0\%A8\%BE\%E0\%A8\%B0}{Punjabi-English Dictionary}, for the weekday; Turner, \href{https://dsal.uchicago.edu/dictionaries/soas/}{CDIAL 11547}, for \emph{vāra}.
\end{tcolorbox}
\section[\texorpdfstring{\pa{ਵੀਰਵਾਰ} (veervaar)}{ਵੀਰਵਾਰ (veervaar)}]{\pa{ਵੀਰਵਾਰ} — Thursday}

\label{lesson:PA-C182-veervar}



\begin{quote}
Say the first three labels: \textbf{\pa{ਸੋਮਵਾਰ}}, \textbf{\pa{ਮੰਗਲਵਾਰ}}, \textbf{\pa{ਬੁੱਧਵਾਰ}}.
\end{quote}

\begin{grammarlens}[title={Grammar Lens: one weekday by ear}]
\emph{Hear:} \emph{veervaar}; YOU SAY: \textbf{\pa{ਵੀਰਵਾਰ}}

\begin{quote}
\textbf{\pa{ਵੀਰਵਾਰ}} — \emph{vīrvār} — Thursday
\end{quote}
\end{grammarlens}

\begin{cousinweb}
\emph{Vīrvār} is a learned borrowing of Sanskrit \emph{vīravāra}. On the surface it contains \emph{vīra}, \textquotedblleft{}hero\textquotedblright{}, and \emph{vāra}, \textquotedblleft{}day\textquotedblright{}.
\end{cousinweb}

\subsection*{Your turn}
\begin{itemize}
\raggedright
  \item \emph{Hear:} \emph{veervaar}; YOU SAY: \textbf{\pa{ਵੀਰਵਾਰ}}
  \item \emph{Read it:} \textbf{\pa{ਵੀਰਵਾਰ}}
  \item \emph{Say it:} \textbf{\pa{ਬੁੱਧਵਾਰ}}, then \textbf{\pa{ਵੀਰਵਾਰ}}
\end{itemize}

\begin{tcolorbox}[breakable,colback=teal!4,colframe=teal!35!black,title={Before you move on}]
Say \emph{Thursday}. (\textbf{\pa{ਵੀਰਵਾਰ}}.)

Sources: Punjabi University, \href{https://dic.learnpunjabi.org/default.aspx?look=\%E0\%A8\%B5\%E0\%A9\%80\%E0\%A8\%B0\%E0\%A8\%B5\%E0\%A8\%BE\%E0\%A8\%B0}{Punjabi-English Dictionary}, for the weekday; \href{https://en.wiktionary.org/wiki/\%E0\%A8\%B5\%E0\%A9\%80\%E0\%A8\%B0\%E0\%A8\%B5\%E0\%A8\%BE\%E0\%A8\%B0}{Wiktionary}, for the learned Sanskrit history.
\end{tcolorbox}
\section[\texorpdfstring{\pa{ਸ਼ੁੱਕਰਵਾਰ} (shukkarvaar)}{ਸ਼ੁੱਕਰਵਾਰ (shukkarvaar)}]{\pa{ਸ਼ੁੱਕਰਵਾਰ} — Friday}

\label{lesson:PA-C182-shukkarvar}



\begin{quote}
Say \textbf{\pa{ਬੁੱਧਵਾਰ}}, then \textbf{\pa{ਵੀਰਵਾਰ}}.
\end{quote}

\begin{grammarlens}[title={Grammar Lens: one weekday by ear}]
\emph{Hear:} \emph{shukkarvaar}; YOU SAY: \textbf{\pa{ਸ਼ੁੱਕਰਵਾਰ}}

\begin{quote}
\textbf{\pa{ਸ਼ੁੱਕਰਵਾਰ}} — \emph{śukkarvār} — Friday
\end{quote}
\end{grammarlens}

\begin{cousinweb}
\emph{Śukkarvār} joins Sanskrit \emph{śukra}, the name of Venus, with \emph{vāra}, \textquotedblleft{}day of the week\textquotedblright{} (CDIAL 11547).
\end{cousinweb}

\subsection*{Your turn}
\begin{itemize}
\raggedright
  \item \emph{Hear:} \emph{shukkarvaar}; YOU SAY: \textbf{\pa{ਸ਼ੁੱਕਰਵਾਰ}}
  \item \emph{Read it:} \textbf{\pa{ਸ਼ੁੱਕਰਵਾਰ}}
  \item \emph{Say it:} \textbf{\pa{ਵੀਰਵਾਰ}}, then \textbf{\pa{ਸ਼ੁੱਕਰਵਾਰ}}
\end{itemize}

\begin{tcolorbox}[breakable,colback=teal!4,colframe=teal!35!black,title={Before you move on}]
Say \emph{Friday}. (\textbf{\pa{ਸ਼ੁੱਕਰਵਾਰ}}.)

Sources: Punjabi University, \href{https://dic.learnpunjabi.org/default.aspx?look=\%E0\%A8\%B8\%E0\%A8\%BC\%E0\%A9\%81\%E0\%A9\%B1\%E0\%A8\%95\%E0\%A8\%B0\%E0\%A8\%B5\%E0\%A8\%BE\%E0\%A8\%B0}{Punjabi-English Dictionary}, for the weekday; Turner, \href{https://dsal.uchicago.edu/dictionaries/soas/}{CDIAL 11547}, for \emph{vāra}.
\end{tcolorbox}
\section[\texorpdfstring{\pa{ਸ਼ਨੀਵਾਰ} (shanivaar)}{ਸ਼ਨੀਵਾਰ (shanivaar)}]{\pa{ਸ਼ਨੀਵਾਰ} — Saturday}

\label{lesson:PA-C182-shanivar}



\begin{quote}
Say \textbf{\pa{ਵੀਰਵਾਰ}}, then \textbf{\pa{ਸ਼ੁੱਕਰਵਾਰ}}.
\end{quote}

\begin{grammarlens}[title={Grammar Lens: one weekday by ear}]
\emph{Hear:} \emph{shanivaar}; YOU SAY: \textbf{\pa{ਸ਼ਨੀਵਾਰ}}

\begin{quote}
\textbf{\pa{ਸ਼ਨੀਵਾਰ}} — \emph{śanīvār} — Saturday
\end{quote}

This is the shorter weekday form used in the draft; other Punjabi materials may use a longer Saturday form.
\end{grammarlens}

\begin{cousinweb}
\emph{Śanīvār} joins Sanskrit \emph{śani}, the name of Saturn, with \emph{vāra}, \textquotedblleft{}day of the week\textquotedblright{} (CDIAL 11547).
\end{cousinweb}

\subsection*{Your turn}
\begin{itemize}
\raggedright
  \item \emph{Hear:} \emph{shanivaar}; YOU SAY: \textbf{\pa{ਸ਼ਨੀਵਾਰ}}
  \item \emph{Read it:} \textbf{\pa{ਸ਼ਨੀਵਾਰ}}
  \item \emph{Say it:} \textbf{\pa{ਸ਼ੁੱਕਰਵਾਰ}}, then \textbf{\pa{ਸ਼ਨੀਵਾਰ}}
\end{itemize}

\begin{tcolorbox}[breakable,colback=teal!4,colframe=teal!35!black,title={Before you move on}]
Say \emph{Saturday}. (\textbf{\pa{ਸ਼ਨੀਵਾਰ}}.)

Sources: Punjabi University, \href{https://dic.learnpunjabi.org/default.aspx?look=\%E0\%A8\%B8\%E0\%A8\%BC\%E0\%A8\%A8\%E0\%A9\%80\%E0\%A8\%B5\%E0\%A8\%BE\%E0\%A8\%B0}{Punjabi-English Dictionary}, for the weekday; IBM, \href{https://public.dhe.ibm.com/software/globalization/documents/India-Punjabi\_Calendar.pdf}{Punjabi locale calendar data}, for the standard weekday set.
\end{tcolorbox}
\section[aitvaar]{Aitvaar — Sunday by ear}

\label{lesson:PA-C182-aitvar}



\begin{quote}
Say \textbf{\pa{ਸ਼ੁੱਕਰਵਾਰ}}, then \textbf{\pa{ਸ਼ਨੀਵਾਰ}}.
\end{quote}

\begin{grammarlens}[title={Grammar Lens: one weekday by ear}]
\emph{Hear:} \emph{aitvaar}; YOU SAY: \emph{aitvaar}

\begin{quote}
\emph{aitvār} — Sunday
\end{quote}

Keep this word oral for now. Its first printed letter comes next.
\end{grammarlens}

\begin{cousinweb}
\emph{Aitvār} is inherited from Sanskrit \emph{ādityavāra}, \textquotedblleft{}Sunday\textquotedblright{}. Its history is older than the transparent weekday compounds beside it.
\end{cousinweb}

\subsection*{Your turn}
\begin{itemize}
\raggedright
  \item \emph{Hear:} \emph{aitvaar}; YOU SAY: \emph{aitvaar}
  \item \emph{Say it:} \textbf{\pa{ਸ਼ਨੀਵਾਰ}}, then \emph{aitvaar}
\end{itemize}

\begin{tcolorbox}[breakable,colback=teal!4,colframe=teal!35!black,title={Before you move on}]
Say \emph{Sunday}. (\emph{Aitvaar}.)

Sources: Punjabi University, Punjabi-English Dictionary, for the weekday; Wiktionary, for its inheritance from Sanskrit. Links are deferred until after the printed form is introduced.
\end{tcolorbox}
\section[\texorpdfstring{\pa{ਐ} (ai)}{ਐ (ai)}]{\pa{ਐ} — one new independent vowel from the Sunday word you can say}

\label{lesson:PA-C182-airhaa-ai}



\begin{quote}
Say \emph{aitvaar}, then \textbf{\pa{ਸ਼ਨੀਵਾਰ}}.
\end{quote}

\subsection*{Script you'll notice: \pa{ਐ}}
\begin{quote}
\pa{ਐ}
\end{quote}

\textbf{aiṛā with dulāvā̃} — \emph{ai}. Keep the independent vowel in view. It opens the Sunday word \emph{aitvaar} that you already know by ear.

No stroke order is taught here: the track does not yet have a verified source for one.

\subsection*{Writing: copy what you can see}
Keep \textbf{\pa{ਐ}} in front of you. Trace its printed shape once with a finger, then copy it once with a pencil while the model stays visible. Say \emph{ai} when you finish.

\begin{quote}
This is observe-and-copy with the model visible, not writing from memory.
\end{quote}

\subsection*{Your turn}
\begin{itemize}
\raggedright
  \item \emph{Look:} at \textbf{\pa{ਐ}}; YOU SAY: \emph{ai}
  \item Look away. Then look back and point to \textbf{\pa{ਐ}}.
  \item \emph{Say it:} \emph{aitvaar}, then point to its opening vowel \textbf{\pa{ਐ}} and say \emph{ai}
  \item \emph{Say it:} \textbf{\pa{ਸ਼ਨੀਵਾਰ}}, then \emph{aitvaar}
\end{itemize}

\begin{tcolorbox}[breakable,colback=teal!4,colframe=teal!35!black,title={Before you move on}]
Point to \textbf{\pa{ਐ}}, say \emph{ai}, then say \emph{aitvaar}.

Source: \href{https://www.unicode.org/charts/PDF/U0A00.pdf}{Unicode Gurmukhi chart}.
\end{tcolorbox}
\section[\texorpdfstring{\pa{ਸੋਮਵਾਰ} \pa{ਤੋਂ} \pa{ਐਤਵਾਰ} (somvaar ton aitvaar)}{ਸੋਮਵਾਰ ਤੋਂ ਐਤਵਾਰ (somvaar ton aitvaar)}]{Monday through Sunday — one guided review}

\label{lesson:PA-C182-weekday-review}



\begin{quote}
Recall the two ends: \textbf{\pa{ਸੋਮਵਾਰ}}, then \textbf{\pa{ਐਤਵਾਰ}}.
\end{quote}

\subsection*{Your turn: rebuild the week}
\begin{quote}
\textbf{\pa{ਸੋਮਵਾਰ}, \pa{ਮੰਗਲਵਾਰ}, \pa{ਬੁੱਧਵਾਰ}, \pa{ਵੀਰਵਾਰ}, \pa{ਸ਼ੁੱਕਰਵਾਰ}, \pa{ਸ਼ਨੀਵਾਰ}, \pa{ਐਤਵਾਰ}}
\end{quote}

Read the line once. Cover it. Rebuild it in two gentle groups: the first four, then the last three.

\begin{itemize}
\raggedright
  \item \emph{Hear:} the seven weekday labels in order; YOU SAY: the same ordered week
  \item \emph{Read it:} \textbf{\pa{ਸੋਮਵਾਰ} \pa{ਤੋਂ} \pa{ਐਤਵਾਰ}}
\end{itemize}

\begin{tcolorbox}[breakable,colback=teal!4,colframe=teal!35!black,title={Before you move on}]
Say the week once without looking. Stop and return to the model if one label drops out.

Source: IBM, \href{https://public.dhe.ibm.com/software/globalization/documents/India-Punjabi\_Calendar.pdf}{Punjabi locale calendar data}, for the ordered weekday set.
\end{tcolorbox}
\section[\texorpdfstring{\pa{ਸੋਮਵਾਰ} \pa{ਤੋਂ} \pa{ਸ਼ੁੱਕਰਵਾਰ} … (somvaar ton …)}{ਸੋਮਵਾਰ ਤੋਂ ਸ਼ੁੱਕਰਵਾਰ … (somvaar ton …)}]{A two-line opening-days notice}

\label{lesson:PA-C182-payoff}



\begin{quote}
Say \textbf{\pa{ਸੋਮਵਾਰ}}, \textbf{\pa{ਸ਼ੁੱਕਰਵਾਰ}}, \textbf{\pa{ਸ਼ਨੀਵਾਰ}}, \textbf{\pa{ਐਤਵਾਰ}}.
\end{quote}

\subsection*{Your turn: four-skill payoff}
First retrieve the whole week once. Then use this small notice:

\begin{quote}
\textbf{\pa{ਸੋਮਵਾਰ} \pa{ਤੋਂ} \pa{ਸ਼ੁੱਕਰਵਾਰ} \pa{ਤੱਕ} \pa{ਖੁੱਲ੍ਹਾ।} \pa{ਸ਼ਨੀਵਾਰ} \pa{ਅਤੇ} \pa{ਐਤਵਾਰ} \pa{ਬੰਦ।}}
\end{quote}

\begin{itemize}
\raggedright
  \item \emph{Hear:} the two-line opening-days notice; YOU SAY: \textbf{\pa{ਸੋਮਵਾਰ} \pa{ਤੋਂ} \pa{ਸ਼ੁੱਕਰਵਾਰ} \pa{ਤੱਕ} \pa{ਖੁੱਲ੍ਹਾ।} \pa{ਸ਼ਨੀਵਾਰ} \pa{ਅਤੇ} \pa{ਐਤਵਾਰ} \pa{ਬੰਦ।}}
  \item \emph{Read it:} \textbf{\pa{ਸੋਮਵਾਰ} \pa{ਤੋਂ} \pa{ਸ਼ੁੱਕਰਵਾਰ} \pa{ਤੱਕ} \pa{ਖੁੱਲ੍ਹਾ।} \pa{ਸ਼ਨੀਵਾਰ} \pa{ਅਤੇ} \pa{ਐਤਵਾਰ} \pa{ਬੰਦ।}}
  \item \emph{Write it:} the same two-line notice after covering the model
\end{itemize}

\begin{tcolorbox}[breakable,colback=teal!4,colframe=teal!35!black,title={Before you move on}]
Rebuild the week, then the notice, without looking.

This is the fifth prerequisite tranche exposed by Punjabi A2 mock drafting. It does not unblock Punjabi A2 mock 1 by itself. It is not an A2 mock, not A2 readiness or attainment, and not external or human validation.

Sources: IBM, \href{https://public.dhe.ibm.com/software/globalization/documents/India-Punjabi\_Calendar.pdf}{Punjabi locale calendar data}, for the weekday set; the already taught Punjabi schedule forms from chapters 178 and 179.
\end{tcolorbox}
